\section{Phân tích và trực quan hóa dữ liệu}

% Slide 9
\begin{frame}[t]{Bức tranh tổng quan 2024}
  \begin{columns}[T,onlytextwidth]
    \begin{column}{0.66\textwidth}
      \begin{figure}
        \centering
        \includegraphics[width=\textwidth,height=0.59\textheight,keepaspectratio]{dashboard/dashboard_overview_page.png}
        \caption*{Bốn góc nhìn về điểm PAPI năm 2024}
      \end{figure}
    \end{column}
    \begin{column}{0.31\textwidth}
      \scriptsize
      \centerline{\textbf{Kỹ thuật trực quan}}
      \begin{itemize}
        \setlength{\itemsep}{3pt}
        \item[-] Bản đồ phân cấp màu.
        \item[-] Biểu đồ thác nước.
        \item[-] Thanh chồng 100\%.
        \item[-] Biểu đồ tần suất.
      \end{itemize}

      \vspace{0.06cm}
      \centerline{\textbf{Kết quả chính}}
      \begin{itemize}
        \setlength{\itemsep}{3pt}
        \item[-] 61 tỉnh có tổng điểm hợp lệ.
        \item[-] Quảng Ninh: $47{,}82$; Kiên Giang: $39{,}91$.
        \item[-] D1-D2 có $r \approx 0{,}72$.
        \item[-] 40/59 tỉnh giảm trong giai đoạn 2018-2024; trung vị $-0{,}94$ điểm.
      \end{itemize}
    \end{column}
  \end{columns}

  \vspace{0.02cm}
  \centering
  {\scriptsize Định hướng: xác định lĩnh vực còn thấp tại từng tỉnh và đối chiếu nhóm tăng điểm để lựa chọn ưu tiên cải thiện.}
\end{frame}
